\subsection*{最初に必要な用語}
\addcontentsline{toc}{subsection}{最初に必要な用語}
\par\smallskip\noindent\refstepcounter{table}\label{tab:ch05-01}\textbf{表 \thetable: 第05章 ソフトウェアテスト 表01}\par\smallskip
\begin{longtable}{P{0.23\linewidth}P{0.69\linewidth}}
\toprule
\textbf{用語} & \textbf{初学者向けの説明} \\
\midrule
テスト対象システム SUT & テストされる対象である。プログラム、アプリケーション、サービス、マイクロサービス、ミドルウェア、システム、システム群、エコシステムまで含み得る。\\
テストケース & 入力値だけでなく、実行条件、時刻条件、手順、期待結果を含む、テスト実行のための具体的な指定である。\\
テストスイート & 複数のテストケースをまとめた集合である。\\
期待結果 & テストを実行したときに、本来得られるべき出力、状態変化、メッセージ、振る舞いである。\\
テストオラクル & 実際の結果が正しいかどうかを判定する基準または仕組みである。人、仕様、モデル、コード注釈、比較ツールなどがなり得る。\\
欠陥 & 障害を引き起こす原因である。コード、設計、要求、設定、運用手順に存在し得る。\\
障害 & 利用者や外部から観測される望ましくない振る舞いである。欠陥が存在しても、条件がそろわなければ障害として現れないことがある。\\
カバレッジ & テストが対象のどの部分をどれだけ通ったかを表す度合いである。文、分岐、条件、要求、状態など、対象は目的によって変わる。\\
回帰テスト & 変更後に、以前できていたことが壊れていないかを再確認するテストである。\\
シフトレフト & テストや品質確認を開発の後半に置くだけでなく、早い段階から行う考え方である。\\
\bottomrule
\end{longtable}
